%%%%%%%%%%%%%%%%%%%- HXN
\def\sode{2}
\begin{name}
	{\tenchude}
	{\tendethi}
	{\tentruong}
	{\thoigian}
\end{name}
\caulc
\Opensolutionfile{ans}[ans/ans-HXN-\sode-T]
%Câu hỏi
\begin{ex}%Câu 1
    Cho hàm số $y=f(x)$ có bảng biến thiên như sau:\\
   \centerline{\begin{tikzpicture}[>=stealth]
           \tkzTabInit[nocadre=false,lgt=1.2,espcl=2.5,deltacl=0.5]{$x$/.7 ,$f'(x)$/.7,$f(x)$/2}
           {$-\infty$ , $-3$ , $0$ , $3$ , $+\infty$}
           \tkzTabLine{ , - , $0$ , + , $0$ , - , $0$ , + , }
           \tkzTabVar{+/$+\infty$ , -/$-1$, +/$1$ , -/$-1$ , +/$+\infty$}
   \end{tikzpicture}}
    Hàm số đã cho đồng biến trên khoảng nào dưới đây?
    \choice
    {$\left(-3;3\right)$}
    {\True $\left(-3;0\right)$}
    {$\left(0;3\right)$}
    {$\left(-\infty;-3\right)$}
\end{ex}

\begin{ex}%Câu 2
    Trong không gian $Oxyz$ , cho hai điểm $A\left(1;1;-2\right)$ và $B\left(2;2;1\right)$. Vectơ $\overrightarrow{AB}$ có tọa độ là
    \choice
    {$\left(-1;-1;-3\right)$}
    {$\left(3;1;1\right)$}
    {\True $\left(1;1;3\right)$}
    {$\left(3;3;-1\right)$}
\end{ex}

\begin{ex}%Câu 3
    Tìm tập xác định của hàm số $y=\log_2\left(x-3\right)$.
    \choice
    {$\mathscr{D}=\left(-\infty;3\right)$}
    {$\mathscr{D}=\mathbb{R}$}
    {\True $\mathscr{D}=\left(3;+\infty\right)$}
    {$\mathscr{D}=\left[3;+\infty\right)$}
\end{ex}

\begin{ex}%Câu 4
    Xác định số hạng đầu $u_1$ và công sai $d$ của cấp số cộng $\left(u_n\right)$ có $u_9=5u_2$ và $u_{13}=2u_6+5$.
    \choice
    {\True $u_1=3$ và $d=4$}
    {$u_1=3$ và $d=5$}
    {$u_1=4$ và $d=5$}
    {$u_1=4$ và $d=3$}
\end{ex}

\begin{ex}%Câu 5
    Họ nguyên hàm của hàm số $f(x)=3x^2+\sin x$ là 
    \choice
    {$x^3+\cos x+C$}
    {$x^3+\sin x+C$}
    {\True $x^3-\cos x+C$}
    {$3x^3-\sin x+C$}
\end{ex}

\begin{ex}%Câu 6
    Cho hình hộp $ABCD.A'B'C'D'$. Vectơ $\vec{v}=\overrightarrow{B'A'}+\overrightarrow{B'C'}+\overrightarrow{B'B}$ bằng vectơ nào dưới đây?
    \choice
    {$\overrightarrow{DB'}$}
    {$\overrightarrow{B'D'}$}
    {$\overrightarrow{BD'}$}
    {\True $\overrightarrow{B'D}$}
\end{ex}

\begin{ex}%Câu 7
    Người ta thống kê khối lượng của $80$ quả măng cụt (đơn vị: gam) và thu được mẫu số liệu sau:\\
    \centerline{\begin{tabular}{|c|c|c|c|c|c|}
            \hline
            Khối lượng (gam) & $\left[80;82\right)$ & $\left[82;84\right)$ & $\left[84;86\right)$ & $\left[86;88\right)$ & $\left[88;90\right)$\\
            \hline
            Số quả & $17$ & $20$ & $25$ & $16$ & $12$\\
            \hline
    \end{tabular}}\\
    Khoảng biến thiên của mẫu số liệu ghép nhóm trên là
    \choice
    {$11$ gam}
    {$12$ gam}
    {\True $10$ gam}
    {$20$ gam}
\end{ex}

\begin{ex}%[2D4N1-2]%[Tổ 19 - Đợt 17 - Chương 4 - Bài 1 - CD - Đề 2]%[Bình]
    Tìm nguyên hàm $F(x)$ của hàm số $ f(x)=3{x^2}+1$, biết $F(1)=3$ là
    \choice
    { ${x^3}+3$}
    { $\dfrac{{x^3}}{3}+x+3$}
    { ${x^3}+x$}
    {\True ${x^3}+x+1$}
    \loigiai
    {
        $\int (3{x^2}+1)\mathrm{\,d}x={x^3}+x+C=F(x)$.\\
        Ta có $F(1)=3\Rightarrow 1+1+C=3\Leftrightarrow C=1$.\\
        Vậy $F(x)={x^3}+x+1$.
    }
\end{ex}

\begin{ex}%Câu 9
    Mỗi ngày bác Mạnh đều đi bộ để rèn luyện sức khỏe. Quãng đường đi bộ mỗi ngày của bác trong 20 ngày được thống kê lại ở bảng sau:\\
    \centerline{\begin{tabular}{|c|c|c|c|c|c|}
            \hline
            Quãng đường & $\left[2,7;3,0\right)$ & $\left[3,0;3,3\right)$ & $\left[3,3;3,6\right)$ & $\left[3,6;3,9\right)$ & $\left[3,9;4,2\right)$\\
            \hline
            Số ngày & 3 & 6 & 5 & 4 & 2\\
            \hline
    \end{tabular}}\\
    Phương sai của mẫu số liệu ghép nhóm gần nhất với giá trị nào sau đây?
    \choice
    {$0{,}19$}
    {$1{,}26$}
    {\True $0{,}13$}
    {$0{,}26$}
\end{ex}

\begin{ex}%Câu 10
    Cho hàm số $y=f(x)$ có đạo hàm $f'(x)=(x-2)(x+1)$, $\forall x\in\mathbb{R}$. Mệnh đề nào dưới đây đúng?
    \choice
    {Hàm số đã cho đồng biến trên $(-1;+\infty)$}
    {Hàm số đã cho nghịch biến trên $(2;+\infty)$}
    {\True Hàm số đã cho nghịch biến trên $(-1;2)$}
    {Hàm số đã cho đồng biến trên $(-1;2)$}
\end{ex}

\begin{ex}%Câu 11
    Trong không gian $Oxyz$ , mặt phẳng đi qua tâm của mặt cầu $\left(x-1\right)^2+\left(y+2\right)^2+z^2=12$ và song song với mặt phẳng $\left(Ox\text{z}\right)$ có phương trình là:
    \choice
    {$y+1=0$}
    {$y-2=0$}
    {\True $y+2=0$}
    {$x+z-1=0$}
\end{ex}

\begin{ex}%Câu 12
    Biết đồ thị $(C)$ của hàm số $y=\dfrac{x^2-4x+5}{x-1}$ có hai điểm cực trị. Đường thẳng đi qua hai điểm cực trị của đồ thị $(C)$ cắt trục hoành tại điểm $M$ có hoành độ $x_M$ bằng
    \choice
    {\True $x_M=2$}
    {$x_M=1-\sqrt{2}$}
    {$x_M=1$}
    {$x_M=1+\sqrt{2}$}
        \end{ex}
\Closesolutionfile{ans}
\cauds
\Opensolutionfile{ans}[ans/ans-HXN-\sode-TF]
%Câu hỏi

\begin{ex}%Câu 13
\immini
{
    Một cái bể nước có dạng khối chóp tứ giác đều ngược với cạnh đáy bằng $3\sqrt{2}$dm và chiều cao bằng $6$dm (tham khảo hình vẽ bên – các kích thước được nêu ra là phần bên trong hình). Nước được bơm vào bể với tốc độ không đổi là $2$ lít/phút và ban đầu bể không chứa nước (các kết quả bên dưới được làm tròn đến hai chữ số thập phân sau dấu phẩy).
    \choiceTF
    {\True Bể nước được bơm đầy sau $18$ phút}
    {\True Tốc độ dâng lên của nước là $0{,}23$ dm/phút khi thể tích nước trong bể bằng $\dfrac{1}{3}$ thể tích của bể}
    {\True Khi mực nước cách miệng bể $0{,}5$ dm, người ta ngừng bơm và bắt đầu xả ra với ước lượng tốc độ giảm chiều cao của mực nước trong bể theo thời gian $t$ (phút) được mô hình hóa bởi hàm số: $h'(t)=\dfrac{1}{350}t-\dfrac{193}{700}$ (dm/phút). Sau $5$ phút, thể tích nước trong bể là $11,97$ dm$^3$}
    {Cùng với dữ kiện của  thì sau $23{,}59$ phút nước trong bể vừa được xả hết}
}
{
    \includegraphics[width=7cm]{img/HXN-2-13}
}
\loigiai{
    \begin{itemchoice}
        \itemch 
        Thể tích chậu nước hình chóp tứ giác đều là $V_{\text{chậu}}=\dfrac{1}{3}\cdot \left(3\sqrt{2}\right)^2\cdot 6=36\,dm^3=36$ lít.\\
        Thời gian bơm nước đầy bể là $\dfrac{36}{2}=18$ phút.
        \itemch Gọi $V$(t), $h$(t) lần lượt là thể tích và chiều cao của nước sau $t$ phút.\\
        Ta có $\dfrac{V(t)}{V_{\text{chậu}}}=\left(\dfrac{h(t)}{6}\right)^3\Rightarrow V(t)=\left(\dfrac{h(t)}{6}\right)^3\cdot 36$ hay $V(t)=\dfrac{(h(t))^3}{6}$\tagEX{*}
        Đạo hàm hai vế của $(*)$ theo biến $t$ ta được: $\dfrac{dV(t)}{\mathrm{\,d}t}=\dfrac{(h(t))^2}{2}\cdot \dfrac{dh(t)}{\mathrm{\,d}t}$\tagEX{**}
        Thời điểm thể tích nước bằng $\dfrac{1}{3}$ thể tích chậu thì \\
        \centerline{
        $\left(\dfrac{h(t)}{6}\right)^3=\dfrac{1}{3}\Rightarrow \dfrac{h(t)}{6}=\sqrt[3]{\dfrac{1}{3}}\Rightarrow h(t)=6\cdot \sqrt[3]{\dfrac{1}{3}}\approx 4{,}16\,dm$.
        }
        Thay $\dfrac{dV(t)}{\mathrm{\,d}t}=2$ lít/phút $=2$ dm/phút; $h(t)\approx 4{,}16$dm vào (**), thì:$\dfrac{dh(t)}{\mathrm{\,d}t}\approx 0{,}23$ dm/phút.
        \itemch Mực nước cách miệng bể $0{,}5$ dm nên chiều cao ban đầu bằng $5{,}5\,dm$.\\
        Chiều cao của nước trong chậu sau $5$ phút là\\
        \centerline{$h(5)=5{,}5+\int\limits_0^5{\left(\dfrac{1}{350}t-\dfrac{193}{700}\right)\mathrm{\,d}t}=\dfrac{291}{70}\approx 4{,}16$dm.}
        Thể tích nước còn lại trong bể là $V(5)=\left(\dfrac{h(5)}{6}\right)^3.36\approx 11{,}97\,dm^3$.
        \itemch Thời gian để mực nước trong chậu đang là $5{,}5$ dm trở về $0$ dm thỏa mãn phương trình\\
        \centerline{$5{,}5+\displaystyle\int\limits_0^t{\left(\dfrac{1}{350}t-\dfrac{193}{700}\right)\mathrm{\,d}t}=0\Rightarrow t\approx 22{,}59$ phút.}
        Với khoảng thời gian $22{,}59$ phút thì nước trong bể vừa được xả hết.
    \end{itemchoice}
}
\end{ex}

\begin{ex}%Câu 14
Có hai tên cướp vừa lấy được một chiếc ca nô ở vị trí $A$ thuộc bờ sông, chúng liền cho ca nô chạy theo phương hợp với bờ sông một góc $60^\circ$ với vận tốc $v=2t$ (mét/giây), trong đó $t$ (giây) là thời gian kể từ khi xuất phát. Sau $21$ giây, ca nô đến vị trí $B$ và chúng quyết định chuyển hướng cho ca nô chuyển động thẳng đều theo phương song song với bờ sông, tầm nửa phút sau thì ca nô đến $C$ (tham khảo hình vẽ).\\
\centerline{
    \includegraphics[width=8cm]{img/HXN-9-14}
}
    \choiceTF
    {\True Vị trí $B$ mà ca nô bọn cướp chuyển hướng cách bờ sông khoảng $382$ m (làm tròn đến hàng đơn vị)}
    {Khoảng cách $A$, $C$ tính theo đường chim bay bằng $1522$ m (làm tròn đến hàng đơn vị)}
    {\True Nếu các chiến sĩ công an khởi động ca nô và đi thẳng từ $D$ đến $C$ với vận tốc được tăng thêm $3$ m sau mỗi giây thì sau $21$ giây sẽ đến vị trí $D$ (làm tròn đến hàng đơn vị của giây)}
    {Trên thực tế các chiến sĩ đã chọn giải pháp cho ca nô khởi động và di chuyển vuông góc với bờ với gia tốc $a$ dương, cùng lúc đó bọn cướp từ vị trí $D$ tiến thẳng về phía trước (giữ nguyên hướng đi và tốc độ), hai bên giáp mặt nhau khi $a=4{,}78$ m/s$^2$ (làm tròn đến hàng phần trăm)}
    \loigiai{
        \begin{center}
            \begin{tikzpicture}[declare function={d=6;b=2;c=3;},thick]
                \path
                (0,0) coordinate (A)
                (-60:b) coordinate (B)
                ($(B)+(0:c)$) coordinate (C)
                (0:d) coordinate (D)
                ($(A)!(B)!(D)$) coordinate (H)
                ($(A)!(C)!(D)$) coordinate (K)
                ($(B)!(D)!(C)$) coordinate (E)
                pic[draw,angle radius=5mm,"$60^\circ$",angle eccentricity=1.7]{angle=B--A--H}
                ;
                \foreach \x/\y/\z in {B/H/K,C/K/D,C/E/D}\draw pic[draw,angle radius =2mm] {right angle = \x--\y--\z};
                \draw[dashed] (B)--(H) (C)--(K) (C)--(E)--(D);
                \draw (A)--(B)--(C) (D)--(A);
                \draw[red] (A)--(C)--(D);
                \foreach \x/\g in {A/180,B/-90,C/-90,D/0,E/-90,H/90,K/90}\draw[fill=white] (\x) circle (1pt)+(\g:3mm) node{$\x$};
            \end{tikzpicture}
        \end{center}
    \begin{itemchoice}
        \itemch Sau $21$ giây, ca nô bọn cướp đi được $AB=\int\limits_0^{21}{2t\mathrm{d}t}=441m$.\\
        Gọi $H$ là hình chiếu vuông góc của $B$ trên bờ sông, khi đó $BH=AB\sin 60^\circ=\dfrac{441\sqrt{3}}{2}\approx 382m$.
        \itemch Vận tốc của ca nô bọn cướp tại $B$ là $v_B=2\times 21=42$ m/s.\\
        Khoảng cách hai vị trí $B$, $C$ là $BC=42\times 30=1260m$.\\
        Khoảng cách hai vị trí $A$, $H$ là $AH=AB\cos 60^\circ=220{,}5m$.\\
        Gọi $K$ là hình chiếu vuông góc của $C$ trên bờ sông thì $HK=BC=1260m$.\\
        Do đó $AK=AH+HK=1480{,}5m$; $CK=BH=\dfrac{441\sqrt{3}}{2}m$ và $AC=\sqrt{AK^2+CK^2}\approx 1529 m$.
        \itemch Ta có $DK=2000-\left(220{,}5+1260\right)=519{,}5m$; suy ra $CD=\sqrt{CK^2+DK^2}\approx 644{,}8m$.\\
        Thời gian để các chiến sĩ đi được từ $D$ đến $C$ thỏa $\displaystyle\int\limits_0^t{3t\mathrm{d}t}=\sqrt{415741}\Rightarrow t\approx 21s$.
        \itemch Gọi $E$ là vị trí hai bên giáp mặt nhau (nếu có) thì tam giác $CDE$ vuông tại $E$.\\
        Khi đó $CE=DK=519{,}5m$ và $DE=\dfrac{441\sqrt{3}}{2}m$.\\
        Thời gian để ca nô bọn cướp đi từ $C$ đến $E$ là $\dfrac{CE}{42}=\dfrac{1039}{84}\approx 12{,}37s$.\\
        Gia tốc $a$ thỏa mãn $\displaystyle\int\limits_0^{\tfrac{1039}{84}}{at \mathrm{d}t}=\dfrac{441\sqrt{3}}{2}\Rightarrow 4{,}99m/s^2$.
    \end{itemchoice}
    }
\end{ex}

\begin{ex}%Câu 15
\immini
{
    Sao Thủy gần như không có khí quyển thật sự như Trái Đất hay sao Kim. Tuy nhiên, nó có một lớp khí rất mỏng gọi là exosphere – tức là thượng quyển loãng, gồm các hạt khí cực kỳ thưa thớt như hydro, heli, oxy, natri...Trong không gian Oxyz, đơn vị trên mỗi trục là nghìn km, vùng thượng quyển loãng của sao Thủy được mô hình hóa bởi phương trình mặt cầu $x^2+y^2+z^2-2x-4y-4=0$. Các nhà khoa học không gian đang quan sát các tiểu hành tinh ở các vị trí có tọa độ $A\left(4;2;4\right)$, $B\left(1;4;2\right)$ và xem xét sự di chuyển của chúng. Nếu tiểu hành tinh nằm trong vùng thượng quyển loãng thì nó sẽ bị hút xuống bề mặt sao Thủy.
}
{
    \includegraphics[width=5cm]{img/HXN-2-15}
}
\choiceTF
    {\True Vùng thượng quyển loãng sao Thủy có tâm $\left(1;2;0\right)$, bán kính bằng $3$}
    {Hai tiểu hành tinh ở các vị trí $A$, $B$ sẽ bị hút xuống bề mặt sao Thủy}
    {Các nhà quan sát cho rằng có một sao chổi mang tên Haxen di chuyển theo quỹ đạo đường thẳng với vận tốc $51{,}5$ km/s; khoảng cách ngắn nhất từ tâm sao Thủy đến sao chổi bằng $\dfrac{\sqrt{871}}{10}$ nghìn km. Thời gian sao chổi đi trong vùng thượng quyển loãng của sao Thủy bằng $20$ giây (làm tròn đến hàng đơn vị của giây)}
    {\True Sao chổi Haxen di chuyển theo phương vectơ $\vec{u}=\left(0;5;2\right)$ . Giả sử $M$, $N$ là điểm đầu và điểm cuối mà sao chổi này đi qua thuộc vùng thượng quyển loãng của sao Thủy. Giá trị nhỏ nhất của tổng $AM+BN$ bằng $3970$ km (làm tròn đến hàng đơn vị của km)}
    \loigiai{
        \begin{itemchoice}
            \itemch Vùng thượng quyển loãng sao Thủy có tâm $I(1;2;0)$, bán kính $R=\sqrt{1^2+2^2+0^2+4}=3$.
            \itemch Ta có $IA=\sqrt{(4-1)^2+(2-2)^2+(4-0)^2}=5$; $IB=\sqrt{(1-1)^2+(4-2)^2+(2-0)^2}=2\sqrt{2}$.\\
            Vì $IA>R$; $IB<R$ nên tiểu hành tinh $A$ nằm ngoài vùng thượng quyển loãng, còn tiểu hành tinh $B$ thì nằm tròn vùng thượng quyển loãng của sao Thủy và nó sẽ bị hút xuống bề mặt sao Thủy.
            \itemch Gọi $d$ là quỹ đạo đường thẳng của sao chổi và $H$ là hình chiếu vuông góc của tâm $I$ trên $d$.\\
            Gọi $M$, $N$ là điểm đầu và điểm cuối của sao chổi trong vùng thượng quyển loãng của sao Thủy.\\
            Ta có $IH=\dfrac{\sqrt{871}}{10}$; suy ra $MN=2MH=2\sqrt{R^2-IH^2}=2\sqrt{9-\dfrac{871}{100}}=\dfrac{\sqrt{29}}{5}$ (nghìn km).
            Thời gian sao chổi di chuyển trong vùng thượng quyển loãng: $\dfrac{\sqrt{29}}{5}\times 1000\colon 51{,}5\approx 21$ (giây).
            \immini
            {
                \itemch Đặt $\vec{MN}=(0;5k;2k)$;\\
                ta có $MN^2=25k^2+4k^2=\dfrac{29}{25}\Rightarrow k=\dfrac{1}{5}$;\\
                suy ra $\vec{MN}=\left(0;1;\dfrac{2}{5}\right)$.\\
            Chọn $C$ sao cho $AMNC$ là hình bình hành;\\
            suy ra $C\left(4;3;\dfrac{22}{5}\right)$.\\
            Khi đó $AM=CN$ và $AM+BN=CN+BN\ge BC=\dfrac{\sqrt{394}}{5}\approx 3{,}97$.\\
            Vậy giá trị nhỏ nhất của biểu thức $AM+BN$ là khoảng $3970$ km.
            }
            {
                \includegraphics[width=4cm]{img/HXN-2-15-LG}
            }
        \end{itemchoice}
    }
\end{ex}

\begin{ex}%Câu 16
\immini
{
    Có hai hộp bóng $A$ và $B$ chỉ đựng các quả bóng đỏ và trắng, trong đó hộp $B$ đựng $4$ quả bóng đỏ và $5$ quả bóng trắng; tổng số bóng hai hộp không qua $20$. Xét hai phép thử ngẫu nhiên sau:\\
Phép thử thứ nhất: Lấy ngẫu nhiên $1$ quả bóng từ hộp $A$ bỏ vào hộp $B$ rồi lấy ngẫu nhiên $1$ quả bóng từ hộp $B$. Bằng thực nghiệm người ta biết được rằng khả năng lấy được quả bóng đỏ từ hộp thứ hai bằng $\dfrac{33}{70}$.\\
Phép thử thứ hai: Lấy ngẫu nhiên $2$ quả bóng từ hộp $A$ bỏ vào hộp B. Sau đó tiếp tục lấy ngẫu nhiên $2$ quả bóng từ hộp $B$.
}
{
\includegraphics[width=5cm]{img/HXN-2-16}
}
\choiceTF
    {Trong phép thử thứ nhất, nếu lấy được quả bóng đỏ từ hộp $A$ bỏ sang hộp $B$ thì xác suất lấy được quả bóng trắng từ hộp $B$ bằng $0{,}4$}
    {Hộp thứ nhất đựng $4$ quả bóng đỏ và $3$ quả bóng trắng
    }
    {Xác suất để lấy được $2$ quả bóng đỏ từ hộp $B$ bằng $\dfrac{166}{1155}$}
    {Nếu biết rằng hai quả bóng lấy được từ hộp $B$ cùng có màu đỏ, xác suất để có ít nhất $1$ quả là từ hộp $A$ chuyển sang bằng $\dfrac{67}{128}$}
    \loigiai{
        \begin{itemchoice}
            \itemch Trong phép thử thứ nhất, nếu lấy được quả bóng đỏ từ hộp $A$ bỏ sang hộp $B$ thì khi đó hộp $B$ có $5$ quả bóng đỏ, $5$ quả bóng trắng. Do đó xác suất lấy được quả bóng trắng từ hộp $B$ bằng $0{,}5$.
            \itemch Trong phép thử thứ nhất, gọi các biến cố $X$: \lq\lq Lấy được quả bóng đỏ từ hộp $A$\rq\rq và $Y$: \lq\lq Lấy được quả bóng đỏ từ hộp thứ hai\rq\rq. Đặt $P(X)=x\in (0;1)$; suy ra $P\left({\bar{X}}\right)=1-x$.
            Ta có sơ đồ hình cây sau đây:\\
            \centerline{
                \includegraphics[width=5cm]{img/HXN-2-16-LG-a}
            }
            Theo giả thiết ta có $P(Y)=0{,}5x+0{,}6(1-x)=\dfrac{33}{70}\Rightarrow x=\dfrac{5}{7}$.\\
            Vì tổng số bóng hai hộp bóng không quá $20$ mà hộp $B$ có $9$ quả bóng nên hộp $A$ có không quá $11$ quả bóng, mà xác suất để lấy được quả bóng đỏ bằng $\dfrac{5}{7}$ nên hộp A có 5 quả bóng đỏ và 2 quả bóng trắng.
            \itemch Trong phép thử thứ hai, ta gọi 2Đ là biến cố \lq\lq Lấy đúng $2$ quả bóng đỏ từ hộp $B$\rq\rq.\\
            Với thông tin có được, ta có sơ đồ hình cây cho phép thử ngẫu nhiên thứ hai:\\
            \centerline{\includegraphics[width=5cm]{img/HXN-2-16-LG-b}}
            Từ đó suy ra $P(2\text{Đ})=\dfrac{\mathrm{C}_5^2}{\mathrm{C}_7^2}\cdot \dfrac{\mathrm{C}_6^2}{\mathrm{C}_{11}^2}+\dfrac{\mathrm{C}_5^1\mathrm{C}_2^1}{\mathrm{C}_7^2}\cdot \dfrac{\mathrm{C}_5^2}{\mathrm{C}_{11}^2}+\dfrac{\mathrm{C}_2^2}{\mathrm{C}_7^2}\cdot \dfrac{\mathrm{C}_4^2}{\mathrm{C}_{11}^2}=\dfrac{256}{1\,155}$.
            \itemch Gọi $X$ là biến cố lấy được ít nhất 1 quả là từ hộp $A$ chuyển sang, ta có\\
            $P\left(X|2\text{Đ}\right)=\dfrac{\dfrac{\mathrm{C}_5^2}{\mathrm{C}_7^2}\cdot \left(\dfrac{2{,}4}{\mathrm{C}_{11}^2}+\dfrac{1}{\mathrm{C}_{11}^2}\right)+\dfrac{5{,}2}{\mathrm{C}_7^2}\cdot \dfrac{1{,}4}{\mathrm{C}_{11}^2}}{P(2\text{Đ})}=\dfrac{65}{128}$.
        \end{itemchoice}
    }
\end{ex}
\Closesolutionfile{ans}
\caukq
\Opensolutionfile{ans}[ans/ans-HXN-\sode-SA]
%Câu hỏi

\begin{ex}%Câu 17
\immini
{
    Từ một khối gỗ hình lập phương có cạnh bằng $5$ dm, người thợ mộc chỉ cần đến hai lát cắt là có thể tạo ra một khối gỗ có dạng hình chóp $S.ABCD$ với đáy $ABCD$ là hình vuông và $SA=AB=5$ dm. Người thợ cần tạo ra một vật để trang trí theo yêu cầu của khách hàng, anh đã chọn $M$ là trung điểm $SB$, $N$ thuộc cạnh $SD$ sao cho $SN=2ND$; sau đó anh ta tiếp tục thực hiện các lát cắt để có được vật thể hình tứ diện $ACMN$, thể tích vật thể sau cùng mà người thợ mộc làm ra là bao nhiêu dm$^3$ (làm tròn đến hàng phần chục)?
\shortans{10,4}
}
{
    \includegraphics[width=5cm]{img/HXN-2-17}
}
\loigiai{
    Gọi $O$ là giao điểm của $AC$ và $BD$ trong mặt phẳng đáy.\\
    \centerline{
        \includegraphics[width=5cm]{img/HXN-2-17-LG}
    }
    Ta có $V_{S.ABCD}=\dfrac{1}{3}SA\cdot S_{ABCD}=\dfrac{5^3}{3}=\dfrac{125}{3}$.\\
    Vì $OM$ là đường trung bình tam giác $SBD$ nên $OM\parallel SD \Rightarrow SD\parallel(AMC)$.\\
    Do đó $d\left(N,(AMC)\right)=d\left(D,(AMC)\right)=d\left(B,(AMC)\right)$.\\
    Suy ra $V_{ACMN}=V_{N.MAC}=V_{D.MAC}=V_{B.MAC}=V_{M.ABC}$.\\
    Ta lại có $\dfrac{V_{M.ABC}}{V_{S.ABCD}}=\dfrac{d\left(M,(ABCD)\right)\cdot S_{ABC}}{d\left(S,(ABCD)\right)\cdot S_{ABCD}}=\dfrac{\dfrac{1}{2}d\left(S,(ABCD)\right)\cdot \dfrac{1}{2}\cdot S_{ABCD}}{d\left(S,(ABCD)\right)\cdot S_{ABCD}}=\dfrac{1}{4}$.\\
    $\Rightarrow V_{M.ABC}=\dfrac{1}{4}V_{S.ABCD}=\dfrac{1}{4}\cdot \dfrac{125}{3}=\dfrac{125}{12}\approx10{,}4$ dm$^3$.
}
    \end{ex}
    
    \begin{ex}%Câu 18
\immini
{
    Jack có một chiếc điện thoại thông minh đã được sạc đầy pin. Nếu Jack không sử dụng điện thoại một phút nào thì máy sẽ hết pin sau $96$ tiếng; còn nếu anh ấy sử dụng điện thoại liên tục thì máy sẽ hết pin sau $8$ tiếng. Biết Jack đã không sử dụng chiếc smartphone trong suốt $36$ tiếng, sau đó lại dùng nó $90$ phút liên tục. Hỏi Jack còn dùng điện thoại được bao nhiêu phút nữa trước khi máy hết pin?
\shortans{210}
}
{
    \includegraphics[width=5cm]{img/HXN-2-18}
}
\loigiai{
    Ta chuẩn hóa tổng thời lượng pin điện thoại ban đầu là $1$.
    \begin{itemize}
        \item Nếu Jack không sử dụng điện thoại thì sau $96$ tiếng smartphone mới hết pin.\\
        Suy ra sau mỗi tiếng không sử dụng, thời lượng pin sẽ giảm đi $\dfrac{1}{96}$.
        \item Nếu Jack dùng điện thoại liên tục thì sau $8$ tiếng, smartphone sẽ hết pin.\\
        Suy ra sau mỗi tiếng sử dụng, thời lượng pin sẽ giảm đi $\dfrac{1}{8}$.
        \item Sau khi Jack không sử dụng điện thoại trong suốt 36 tiếng, thời lượng pin giảm $36\times \dfrac{1}{96}=\dfrac{3}{8}$.\\
        Thời lượng pin còn lại là: $1-\dfrac{3}{8}=\dfrac{5}{8}$.
        \item Sau khi Jack sử dụng điện thoại liên tục trong $90$ phút, tức $\dfrac{3}{2}$ tiếng, thời lượng pin tiếp tục giảm đi: $\dfrac{3}{2}\times \dfrac{1}{8}=\dfrac{3}{16}$.\\
        Thời lượng pin còn lại là $\dfrac{5}{8}-\dfrac{3}{16}=\dfrac{7}{16}$.
    \end{itemize}
    Vậy trước khi điện thoại hết pin, Jack còn có thể sử dụng $\dfrac{7}{16}: \dfrac{1}{8}=3{,}5$ giờ = $210$ phút. 
}
\end{ex}

\begin{ex}%Câu 19
\immini
{
    Một quả trứng khủng long đồ chơi bằng nhựa có thiết diện qua trục lớn là một đường elip. Biết độ dài mỗi trục là $12$ cm và $8$ cm.\\
Bên trong quả trứng người ta cần thiết kế một chiếc hộp hình trụ để đựng các đồ chơi trẻ con như bóng đèn xanh đỏ, kẹo v.v...\\
Hỏi khối trụ như thế có thể tích tối đa bao nhiêu cm$^3$ (làm tròn đến hàng đơn vị).
\shortans{232}
}
{
    \includegraphics[height=5cm]{img/HXN-2-19-a}\includegraphics[width=5cm]{img/HXN-2-19-b}
}
\loigiai{
    Gắn elip lên hệ trục tọa độ $Oxy$ như hình vẽ, elip có $2a=12\Rightarrow a=6$; $2b=8\Rightarrow b=4$.\\
    \centerline{
    \includegraphics[width=5cm]{img/HXN-2-19-LG}
    }
    Phương trình chính tắc elip $(E)$: $\dfrac{x^2}{36}+\dfrac{y^2}{16}=1$.\\
    Đặt chiều cao và bán kính đáy hình trụ nội tiếp elip là $h$, $r$ thì điểm tiếp xúc  $M\left(\dfrac{h}{2};r\right)$ với $0<h<12;0<r<4$.\\
    Điểm $M$ thuộc $(E)$: $\dfrac{x^2}{36}+\dfrac{y^2}{16}=1\Rightarrow \dfrac{\left(\dfrac{h}{2}\right)^2}{36}+\dfrac{r^2}{16}=1 \Rightarrow \dfrac{h^2}{144}+\dfrac{r^2}{16}=1\Rightarrow r^2=16\left(1-\dfrac{h^2}{144}\right)$.\\
    Thể tích khối trụ là $V=\pi r^2h=\pi 16\left(1-\dfrac{h^2}{144}\right)\cdot h=\pi \left(16h-\dfrac{h^3}{9}\right)$ hay $V=\pi \left(16h-\dfrac{h^3}{9}\right)$.\\
    Ta có $V'=\pi \left(16-\dfrac{h^2}{3}\right)$; $V'=0\Rightarrow 16-\dfrac{h^2}{3}=0\Rightarrow h^2=48\Rightarrow h=4\sqrt{3}\in (0;12)$.\\
    Giá trị lớn nhất của thể tích khối trụ là $V_{\max }=V\left(4\sqrt{3}\right)=\pi \left(16.4\sqrt{3}-\dfrac{\left(4\sqrt{3}\right)^3}{9}\right)\approx 232$ cm$^3$.
}
\end{ex}

\begin{ex}%Câu 20
Một nhóm học sinh lớp $12$ đã lên bản thiết kế mẫu hoa văn cho một loại gạch men lát nền nhà. Các em đã vẽ $4$ đường cong như hình, từ đó tạo thành một hình $(\mathscr{H})$ khép kín ở giữa viên gạch để tạo điểm nhấn. Cụ thể cách dựng hình được thực hiện như sau:
\begin{itemize}
    \item Dựng hệ trục $Oxy$ với điểm $O$ là tâm của viên gạch, tia $Ox$ hướng sang phải và tia $Oy$ hướng lên trên, đơn vị trên mỗi trục là $5$ cm
    \item Các em lấy $O$ là tâm viên gạch và $A$ là trung điểm một cạnh viên gạch, xác định được điểm $B$ thỏa mãn $\overrightarrow{OB}=\dfrac{5}{6}\overrightarrow{OA}$.
    \item Dựng đường thẳng $\Delta \colon 5x-9=0$. Đường cong $\left(L_1\right)$ là tập hợp các điểm $M$ thỏa mãn $3MB=5d\left(M,\Delta\right)$.
    \item Lấy đối xứng đường cong $\left(L_1\right)$ qua tâm $O$ và qua các đường chéo của viên gạch thì được các đường cong còn lại.
\end{itemize}
\begin{center}
    \includegraphics[width=5cm]{img/HXN-2-20-a}\qquad  \includegraphics[width=5cm]{img/HXN-2-20-b}
\end{center}
Biết viên gạch là hình vuông có kích thước $60$cm; hỏi diện tích hình $(\mathscr{H})$ là bao nhiêu centimét vuông (làm tròn đến hàng đơn vị)?
\shortans{1168}
\loigiai{
    Ta chọn hệ trục tọa độ như hình vẽ với mỗi đơn vị trên trục bằng $5$ cm.\\
    \centerline{
        \includegraphics[width=5cm]{img/HXN-2-20-LG}
    }
    Đường cong $\left(L_1\right)$ là tập hợp điểm $M$ thỏa $3MB=5d\left(M,\triangle \right)$ hay $\dfrac{MB}{d\left(M,\triangle \right)}=\dfrac{5}{3}$\tagEX{1}
    Phương trình $\triangle \colon x=\dfrac{9}{5}$ \tagEX{2}
    Từ $(1)$ và $(2)$ ta thấy $M$ thuộc một nhánh của hyperbol $\left(L_1\right)$: $\dfrac{x^2}{a^2}-\dfrac{y^2}{b^2}=1$; trong đó điểm $B(5;0)$ là một trong hai tiểu điểm của $\left(L_1\right)$ nên $c=5$.\\
    Đường chuẩn $x=\dfrac{a}{e}=\dfrac{a^2}{c}=\dfrac{9}{5}\Rightarrow a=3$ (thử lại ta thấy $\dfrac{MB}{d\left(M,\triangle \right)}=\dfrac{5}{3}=e$ (hợp lí)).\\
    Do $c^2=a^2+b^2\Rightarrow b=\sqrt{c^2-a^2}=4$.\\
    Phương trình $\left(L_1\right)\colon \dfrac{x^2}{9}-\dfrac{y^2}{16}=1$ $(x>3)$.\\
    Xét giao điểm của $\left(L_1\right)$ với đương thẳng $y=x$, ta có $\dfrac{x^2}{9}-\dfrac{x^2}{16}=1\,\,\left(x>3\right)\Rightarrow x=\dfrac{12\sqrt{7}}{7}\approx 4{,}54$.\\
    Từ phương trình $\dfrac{x^2}{9}-\dfrac{y^2}{16}=1\Rightarrow \dfrac{y^2}{16}=\dfrac{x^2}{9}-1\Rightarrow y=4\sqrt{\dfrac{x^2}{9}-1}$ $(x>3)$.\\
    Diện tích cần tính là $S=8\left(\int\limits_0^3{x\mathrm{\,d}x}+\int\limits_3^{\tfrac{12\sqrt{7}}{7}}{\left(x-4\sqrt{\dfrac{x^2}{9}-1}\right)\mathrm{\,d}x}\right)\times 25\approx 1168$ cm$^2$.
}
\end{ex}

\begin{ex}%Câu 21
Trong không gian $Oxyz$, cho ba điểm $A\left(-8;-1;6\right)$, $B\left(1;2;3\right)$, $C\left(-4;14;\sqrt{11}\right)$. Điểm $M$ di động trên mặt cầu $\left(S_1\right)\colon {(x-4)^2}+(y-3)^2+(z+3)^2=49$ sao cho tam giác $MAB$ có $2\sin\widehat{MAB}=\sin\widehat{MBA}$. Tính giá trị nhỏ nhất của đoạn thẳng $CM^2$ (làm tròn đến hàng đơn vị).
\shortans{64}
\loigiai{
    Xét $\triangle MAB$, ta có $\dfrac{BM}{\sin \widehat{MAB}}=\dfrac{AM}{\sin \widehat{MBA}}=2R\Rightarrow \sin \widehat{MAB}=\dfrac{BM}{2R};\sin \widehat{MBA}=\dfrac{AM}{2R}$.\\
    Theo giả thiết $2\sin\widehat{MAB}=\sin\widehat{MBA} \Rightarrow 2\cdot \dfrac{BM}{2R}=\dfrac{AM}{2R}\Rightarrow AM=2BM$.\\
    Gọi $M(x;y;z)$ thì ta có 
    \allowdisplaybreaks
    \begin{eqnarray*}
        &&AM^2=4BM^2\\
        &\Leftrightarrow& (x+8)^2+(y+1)^2+(z-6)^2=4\left[(x-1)^2+(y-2)^2+(z-3)^2\right]\\ &\Leftrightarrow& 3x^2+3y^2+3z^2-24x-18y-12z-45=0\\
        &\Leftrightarrow& x^2+y^2+z^2-8x-6y-4z-15=0
    \end{eqnarray*}
    Do đó $M$ di động trên mặt cầu $(S_2)$ có tâm $I_2(4;3;2)$, bán kính $R_2=2\sqrt{11}$.\\
    Mặt khác ta cũng có $M\in (S_1)$ có tâm $I_1(4;3;-3)$, bán kính $R_1=7$.\\
    Ta có $I_1I_2=5<R_1+R_2$ nên $M$ thuộc đường tròn $(C)=(S_1)\cap (S_2)$.
    \immini
    {
        Tập hợp điểm $M$ thuộc $(C)$ thỏa hệ phương trình\\
    $\heva{& (x-4)^2+(y-3)^2+(z+3)^2=49 \\& x^2+y^2+z^2-8x-6y-4z-15=0} $\\
    $\Leftrightarrow \heva{& x^2+y^2+z^2-8x-6y+6z-15=0 \\& x^2+y^2+z^2-8x-6y-4z-15=0 }\\ \Leftrightarrow \heva{& z=0\,\,(Oxy) \\& x^2+y^2+z^2-8x-6y-4z-15=0 } $.
    }
    {
        \includegraphics[width=5cm]{img/HXN-2-21-LG}
    }
    Vậy $(C)$ thuộc mặt phẳng $(Oxy)$; hình chiếu của $I_1(4;3;-3)$ trên $(Oxy)$ là điểm $E(4;3;0)$ cũng là tâm của đường tròn $(C)$; bán kính $(C)$ là $r=\sqrt{R_1^2-I_1E^2}=\sqrt{7^2-9}=2\sqrt{10}$.\\
    Gọi $C'(-4;14;0)$ là hình chiếu của $C$ lên mặt phẳng $(Oxy)$, ta có $EC'=\sqrt{185}>r$ nên $C'$ nằm ngoài đường tròn $(E; r)$.\\
    Ta có $CM^2\ge CM_0^2=CC'^2+C'M_0^2=CC'^2+\left(EC'-r\right)^2 =11+\left(\sqrt{185}-2\sqrt{10}\right)^2\approx 64$.\\
    Vậy giá trị nhỏ nhất của $CM^2$ xấp xỉ $64$.
}
\end{ex}

\begin{ex}%Câu 22
\immini
{
    Vào một hội thi thiết kế đèn lồng Trung thu, ban tổ chức nhận được một chiếc đèn lồng đặc biệt có hình một tứ diện đều. Trên mỗi cạnh tứ diện thí sinh thiết kế $3$ bóng đèn nằm ở $3$ vị trí chia cạnh tứ diện thành $4$ đoạn bằng nhau. Cứ mỗi phút trôi qua, sẽ có ngẫu nhiên $3$ bóng đèn phát sáng, các bóng còn lại thì tắt. Tính xác suất để ngay phút đầu tiên được ban giám khảo chấm điểm, có $3$ bóng đèn phát sáng ứng với $3$ điểm tạo nên mặt phẳng song song với đúng một cạnh của tứ diện, biết rằng $3$ bóng đèn không hoàn toàn thuộc về một cạnh tứ diện. (Kết quả được làm tròn đến hàng phần trăm).
\shortans{0,27}
}
{
     \includegraphics[width=5cm]{img/HXN-2-22}
}
\loigiai{
    Tổng số cách chọn ra $3$ trong $18$ điểm là $\mathrm{C}_{18}^3=816$.\\
    Tuy nhiên, sẽ có các trường hợp ba điểm thẳng hàng, đó là khi ta lấy ba điểm thuộc cùng một cạnh, tổng số cách là $6\times \mathrm{C}_3^3=6$.\\
    Do đó $n\left(\Omega \right)=816-6=810$.\\
    Xét các mặt phẳng qua $3$ điểm ($3$ bóng đèn) và song song với đoạn $AB$.
    \immini
    {
        Trường hợp $1$: Chọn $1$ cặp điểm thuộc mặt phẳng $(ABC)$ và $1$ điểm không thuộc $(ABC)$.
    \begin{itemize}
        \item Bước $1$: Có $3$ cách chọn $1$ cặp điểm thuộc $(ABC)$.
        \item Bước $2$: Với mỗi cách chọn trong bước $1$ thì lẽ ra sẽ có $9$ cách chọn $1$ điểm không thuộc mặt phẳng $(ABC)$; tuy nhiên vì điều kiện mặt phẳng qua $3$ điểm chỉ song song đúng $1$ cạnh tứ diện (ở đây là $AB$) nên ta loại $3$ điểm trong số $9$ điểm không thuộc $(ABC)$ (ví dụ khi chọn cặp điểm $M$, $N$ thuộc $(ABC)$ thì ta loại $P$, $Q$, $R$  thuộc $AD$, $BD$, $CD$).\\
        Do vậy ta có $3\times 6=18$ cách chọn bộ ba điểm trong trường hợp này.
    \end{itemize}
    Trường hợp $2$: Chọn $1$ cặp điểm thuộc mặt phẳng $(ABD)$ và $1$ điểm không thuộc $(ABD)$.
    }
    {
        \includegraphics[width=5cm]{img/HXN-2-22-LG}
    }
    Ta cũng có $3\times 6=18$ cách chọn bộ ba điểm thỏa mãn.\\
    Vì tính chất bình đẳng của $6$ cạnh trong tứ diện đều, ta có tất cả $6\cdot (18+18)=216$ bộ ba điểm thỏa mãn.\\
    Vì vậy xác suất cần tính là $P=\dfrac{216}{810}=\dfrac{4}{15}\approx 0{,}27$.
}
\end{ex}
\Closesolutionfile{ans}
\inputansbox{6,4,3}{ans/ans-HXN-\sode-T,ans/ans-HXN-\sode-TF,ans/ans-HXN-\sode-SA}